An integer has exactly three divisors precisely when it is $p^2$ for a prime $p$: its divisors are $1,p,p^2$, and the divisor-count formula allows no other prime-exponent pattern. Here $p\le100$. The primes in that range are
\[2,3,5,7,11,13,17,19,23,29,31,37,41,43,47,53,59,61,67,71,73,79,83,89,97.\]
There are $\boxed{25}$ of them.
